%!TEX root = ../main.tex
% =============================================================================
%  图 2  四问的模型依赖关系（正文 §1.3）
%  规格：附件 -> Q1 -> Q2 -> Q3 -> Q4，每个箭头必须标注"提供什么"；
%        本图用于证明四问是一条信息结构递进链，而不是四篇互不相干的报告。
% =============================================================================
\begin{tikzpicture}[
  font=\small, >=Stealth,
  att/.style ={draw=gray!50, fill=gray!8, rounded corners=2pt, align=center,
               font=\scriptsize, text width=2.0cm, minimum height=0.72cm, inner sep=3pt},
  pre/.style ={draw=orange!55!black, fill=orange!8, rounded corners=2pt, align=center,
               font=\scriptsize, text width=6.6cm, minimum height=0.85cm, inner sep=4pt},
  qbox/.style={draw=blue!50, fill=blue!7, rounded corners=2pt, align=center,
               font=\small, text width=4.6cm, minimum height=0.95cm, inner sep=4pt},
  obox/.style={draw=green!45!black, fill=green!7, rounded corners=2pt, align=center,
               font=\scriptsize, text width=2.2cm, minimum height=0.8cm, inner sep=3pt},
  ar/.style  ={->, draw=gray!65, line width=0.6pt},
  lab/.style ={font=\scriptsize, anchor=west, align=left, text width=5.6cm,
               inner sep=1pt, gray!20!black},
]

% ---------------- 附件 ----------------
\node[att] (a1) at (-4.35,0) {附件 1\\电价/负载/光伏};
\node[att] (a2) at (-1.45,0) {附件 2\\负载/光伏实际值};
\node[att] (a3) at ( 1.45,0) {附件 3\\四个时刻光伏预报};
\node[att] (a4) at ( 4.35,0) {附件 4\\波动电价};

% ---------------- 数据预处理 ----------------
\node[pre] (pre) at (0,-1.35)
  {数据预处理：时标对齐（标签 $=$ 时段末端）\ \textbullet\ 附件 1 身份核查\ \textbullet\ 特征与误差分析（第三章）};

\foreach \n in {a1,a2,a3,a4}{\draw[ar] (\n.south) -- (pre.north);}

% ---------------- 四问 ----------------
\node[qbox] (q1) at (0,-2.95) {\textbf{问题 1}\ 确定性线性规划（LP）};
\node[qbox] (q2) at (0,-4.60) {\textbf{问题 2}\ 两阶段随机规划};
\node[qbox] (q3) at (0,-6.25) {\textbf{问题 3}\ 多阶段滚动 $+$ 非预期性};
\node[qbox] (q4) at (0,-7.90) {\textbf{问题 4}\ 波动电价 $+$ 风险度量};

\draw[ar] (pre.south) -- (q1.north);

% 箭头上的"提供什么"标注
\draw[ar] (q1.south) -- (q2.north);
\node[lab] at (0.35,-3.78)
  {\textbf{提供}\ 功率平衡与储电量递推骨架；无储能基线 $C^{base}$ 口径；
   储电水价值与阈值律语言};

\draw[ar] (q2.south) -- (q3.north);
\node[lab] at (0.35,-5.43)
  {\textbf{提供}\ 固定场景集（消融与 VoI 共用）；报童临界分位数给出的安全边际；
   紧急购电的结构性认识};

\draw[ar] (q3.south) -- (q4.north);
\node[lab] at (0.35,-7.08)
  {\textbf{提供}\ 多阶段滚动框架与非预期性实现；终端条件处理；
   预报频次对照（VoI）方法};

% ---------------- 输出 ----------------
\node[obox] (o1) at (-5.55,-2.95) {result1.xlsx\\+ 表 1 / 表 2};
\node[obox] (o2) at (-5.55,-4.60) {result2.xlsx\\+ 表 3};
\node[obox] (o3) at (-5.55,-6.25) {result3.xlsx\\+ VoI 结论};
\node[obox] (o4) at (-5.55,-7.90) {result4-2.xlsx\\result4-3.xlsx};

\foreach \q/\o in {q1/o1,q2/o2,q3/o3,q4/o4}{
  \draw[ar, draw=green!45!black] (\q.west) -- (\o.east);
}

\end{tikzpicture}
